%=============================================================================!
% 16.5  HOW FAR ATOMS WANDER                                                  !
%=============================================================================!
\section{How far atoms wander}
\label{sec:ob-msd}

A drop of ink in a glass of still water spreads without any current to
carry it: each grain of dye is knocked about by the molecules around it,
and the cloud of grains widens. No grain goes anywhere in particular, yet
the width of the cloud grows steadily with time. Atoms in a liquid wander
in the same way, and the rate at which the width grows is the diffusion
coefficient, which Chapter~13 measured as a sixth of a slope and promised
to treat properly here.

\subsection*{A random walk}

A walker who takes steps $\vb{s}_1, \vb{s}_2, \dots$, each of random
direction and independent of the others, is at $\vb{r}_n = \vb{s}_1 +
\dots + \vb{s}_n$ after $n$ of them. The square of that sum is the sum of
all products of two steps, so
\begin{equation*}
   \ens{\lvert\vb{r}_n\rvert^2}
   = \sum_{k=1}^n\ens{\lvert\vb{s}_k\rvert^2}
   + \sum_{k\neq l}\ens{\vb{s}_k\cdot\vb{s}_l}
   = n\ens{s^2} ,
\end{equation*}
the cross terms vanishing because independent steps of mean zero have
zero covariance (\cref{sec:cv-correlation}). The mean squared distance
from the start grows in proportion to the number of steps, and so to the
time, while the mean position stays where it began. The diffusion
coefficient $D$ is defined by that growth: for an atom in $d$ directions,
\begin{equation}
   \ens{\lvert\vb{r}(t_0 + t) - \vb{r}(t_0)\rvert^2} \to 2dDt
   \qquad \text{as $t$ grows},
   \label{eq:ob-einstein}
\end{equation}
which is $6Dt$ in three dimensions, the six of Chapter~13. The left side is
the mean squared displacement, the MSD, averaged over atoms and over
starting times $t_0$.

The factor $2d$ makes $D$ the coefficient of Fick's law, by which a
difference of density drives a flow of atoms. In one dimension let
$\rho(x, t)$ be the number density of some marked atoms and let them
flow at the rate $j = -D\,\partial\rho/\partial x$ per unit area and
time, down the gradient. Atoms are neither made nor lost, so the number
between $x$ and $x + \dd x$, $\rho\,\dd x$ per unit area, changes by
what flows in at $x$ less what flows out at $x + \dd x$, $j(x) - j(x +
\dd x) = -(\partial j/\partial x)\,\dd x$, which gives $\partial\rho/\partial
t = -\partial j/\partial x = D\,\partial^2\rho/\partial x^2$. If the marked
atoms all start at $x = 0$, their mean square distance $\sigma_x^2 = \int
x^2\rho\,\dd x\big/\int\rho\,\dd x$ changes as
\begin{equation*}
   \frac{\dd\sigma_x^2}{\dd t}
   = \frac{D\int x^2\,\partial^2\rho/\partial x^2\,\dd x}{\int\rho\,\dd x}
   = \frac{D\int 2\rho\,\dd x}{\int\rho\,\dd x} = 2D ,
\end{equation*}
the denominator being constant because atoms are kept. The middle step
integrates by parts twice (\cref{sec:la-euler}), with $\rho$ and its slope
vanishing far away: $\int x^2\rho''\,\dd x = -\int 2x\rho'\,\dd x = \int
2\rho\,\dd x$, primes marking slopes in $x$. So $\sigma_x^2 = 2Dt$ in each
direction, and the $d$ directions add.

\subsection*{The mean squared displacement of a run}

Three things make \cref{eq:ob-einstein} work on a trajectory. The
positions must be unwrapped: an atom that leaves the cell and is put back
at the far side has not jumped a cell's width, and the integrator of
\texttt{md.run} keeps each path continuous for this reason
(\cref{sec:md-periodic}). The centre of mass must be taken out of every
frame, so that a drift of the whole system does not count as wandering.
A run started with velocities whose total momentum was not removed drifts
for ever, and a thermostat that does not conserve momentum lets the
centre wander: Langevin's friction and random kicks act on each atom
separately, so their sums over the atoms do not cancel, and the total
momentum performs a random walk of its own. And the average should run over every starting time, not
the first frame alone, since every stretch of a settled run is as good a
start as any other. The function \texttt{msd} of
\texttt{mdlab.analysis.transport} does all three, returning the MSD along
each axis; \texttt{diffusion\_coefficient} fits a straight line to a
window of it by least squares (\cref{sec:cv-drift}) and divides the slope
by $2d$.

At short times the MSD is not linear. Over a time short against that
between collisions an atom keeps its velocity, so its displacement is
$\vb{v}t$ and $\ens{\lvert\vb{v}t\rvert^2} = (3\kB T/m)\,t^2$ by
equipartition (\cref{sec:sm-equipartition}): the motion is ballistic, the
MSD growing as $t^2$. The two forms meet where $(3\kB T/m)t^2 = 6Dt$, at
$t = 2Dm/\kB T$, which for liquid argon is \SI{274}{\femto\second}.
\Cref{fig:ob-msd}(a) shows both regimes for the argon of Chapter~15 at
fixed energy, from a run of \SI{100}{\pico\second} whose paths are summed
from velocities saved every step. After a picosecond the MSD is a
straight line of slope $6D$, and the fit over 2 to \SI{20}{\pico\second}
gives $D = \SI{3.763e-4}{\angstrom\squared\per\femto\second}$ at
\SI{132.1}{\kelvin}. A fit must start beyond the ballistic part: one over
0.02 to \SI{0.2}{\pico\second} gives a third less, while windows from
$t_1$ to $10t_1$ with $t_1$ from \SI{0.05}{\pico\second} on agree within
6\% (\cref{sec:ob-trajectory}).
The ratio MSD$/6t$ approaches $D$ more slowly, reaching 0.9 of it only at
\SI{0.5}{\pico\second}, because the straight line does not pass through
the origin.

The error comes from blocks of the run, as in Chapter~15. The
\SI{2}{\nano\second} run at fixed energy of \cref{sec:cv-observables} cut
into 20 blocks of \SI{100}{\pico\second} gives $D = (3.834 \pm
0.034)\times10^{-4}$\,\si{\angstrom\squared\per\femto\second} at
\SI{133.70}{\kelvin}, its value there. Fits that use only the first frame
of each block as the start spread 2.6 times as widely as those that use
every frame, which is the reason to average over starting times. The
value depends on the size of the cell as $1/L$ (\cref{sec:cv-size}).

\paragraph{Checked against ASE.} ASE's \texttt{DiffusionCoefficient}
fits each axis's MSD, from the first frame of each segment alone and
with no centre removed, over every lag. On four segments of
\SI{50}{\pico\second} of the same run, \texttt{msd} with one start per
segment and no centre removed gives the same slopes to a relative
\num{1e-15}. The ASE estimate is the noisier of the two for the reason
above; the comparison checks the arithmetic, and ASE's time unit, which
is $1/\sqrt{\texttt{FORCE\_TO\_ACCEL}}$ femtoseconds.

\subsection*{Water}

Chapter~11 left the oxygen atoms of its rigid water wandering by
\num{1.41}, 4.14 and \SI{5.76}{\angstrom}, root mean square from the first
frame, after 1, 5 and \SI{10}{\pico\second} and promised a diffusion coefficient. Its
\SI{10}{\pico\second} run of 64 molecules at \SI{304}{\kelvin}, fitted over
1 to \SI{5}{\pico\second}, gives $D = (5.24 \pm 0.47)\times10^{-4}$\,\si{
\angstrom\squared\per\femto\second}, the error from the scatter of the 64
molecules' own fits. A run so short gives $D$ to a tenth; the 64-molecule
box would also need the correction for its size of \cref{sec:cv-size}.

\begin{figure}[htb]
\centering
\includegraphics{ch16_observables/msd}
\caption{Liquid argon at fixed energy and \SI{132.1}{\kelvin}
(\SI{100}{\pico\second}, paths summed from velocities saved every
\SI{10}{\femto\second}). (a) The mean squared displacement against time,
with the ballistic $3\kB Tt^2/m$ (dashed) and $6Dt$ (dotted). (b) Three
estimates of $D$ against the time they reach: MSD$/6t$, the slope of the
MSD divided by 6 (thick, pale), and the Green-Kubo integral
$\frac13\int_0^tC_{vv}$ of \cref{sec:ob-vacf} (thin), which coincide; the
dotted line is $D$ from the fit over 2 to \SI{20}{\pico\second}.}
\label{fig:ob-msd}
\end{figure}

\begin{keyidea}
The mean squared displacement grows as $2dDt$ once the motion forgets
its velocity, and as $(3\kB T/m)t^2$ before. $D$ is fitted to its linear
part, from unwrapped paths with the centre of mass removed and averaged
over every starting frame, with its error from blocks and a correction
for the size of the cell.
\end{keyidea}

\notebook{16}{move the fitting window across the MSD of liquid argon}

\begin{selfcheck}
An atom of mass \SI{40}{\amu} in a liquid at \SI{300}{\kelvin} has $D =
\SI{3e-4}{\angstrom\squared\per\femto\second}$. Roughly when does its
motion stop being ballistic?
\end{selfcheck}
